14.0 SKARBY

Ogólna zasada:

Jeżeli oddział pokona w walce potwora lub grupę potworów (a także czasem po zakończonych sukcesem negocjacjach) gracze mogą ustalić jakie skarby potwór miał ze sobą lub też jakich skarbów strzegł.
Skarby dzielą się na trzy rodzaje: sztuki złota, kosztowności oraz magiczne przedmioty.

Po dokładnym ustaleniu rodzaju i ilości skarby dzielone są między członków oddziału w sposób, jaki gracze uznają za właściwy.
Skarby, posiadane przez każdą postać muszą znaleźć swoje odbicie na karcie cech tej postaci.
Postać, która otrzymała magiczny przedmiot może go używać do końca gry, odstąpić, wymienić lub odsprzedać.

Jeżeli jakaś postać zostanie zabita pozostali członkowie oddziału mogą podzielić między siebie przez nią skarby.

Sposób postępowania:

1.
Z tabeli ,,Charakterystyka potworów'' należy odczytać typ skarbu, odpowiadający zabitemu potworowi.
Jest on określany jedną lub dwoma literami, przedzielonymi kreską.
W przypadku dwóch liter – pierwsza z nich oznacza typ skarbu, należący do potwora osiadłego, zaś druga jest używana wtedy, gdy zabity potwór był potworem wędrownym.

2.
Dokładną liczbę poszczególnych rodzajów skarbów określa się przy pomocy tabeli 14.9 ,,Skarby''.
Ta tabela posiada trzy kolumny, odpowiadające poszczególnym rodzajom skarbów.
Każdemu typowi skarbu (tzn. literze) odpowiada kod, który odczytywany jest w następujący sposób: pierwsza cyfra (np. 3:1K6x5) oznacza prawdopodobieństwo, że potwór w ogóle posiada ten rodzaj skarbu.
Należy rzucić 1K6 i jeżeli liczba wyrzuconych oczek jest mniejsza od tej cyfry lub jej równa, to znaczy, że potwór posiada ten rodzaj skarbu (w przykładzie kodu podanym wyżej – 1, 2 lub 3 oczka).

Z drugiej strony dwukropka znajduje się kodowe oznaczenie kostki (jednej lub więcej tego samego rodzaju), którą należy rzucić, aby dowiedzieć się jakiej wielkości jej skarb – w liczbach sztuk złota, ilości kosztowności (ilość, a nie wartość) lub liczby magicznych przedmiotów równa się liczbie wyrzuconych oczek.
W przykładzie kodu podanego wyżej – należy jeszcze raz pomnożyć przez pięć.

Jeżeli potwór posiadał kosztowności należy ustalić ich wartość.
Służy do tego tabela 14.9 ,,Kosztowności''.
Gracz musi rzucić 2K6 osobno dla każdej odebranej potworowi kosztowności i odczytać z tabeli jej wartość, którą znajduje się w rzędzie, odpowiadającym liczbie wyrzuconych oczek.
Wartość kosztow-

ności powinna być zapisywana na karcie cech postaci.

3.
Jeżeli potwór posiadał magiczne przedmioty należy ustalić ich rodzaj przy pomocy tabeli 14.9 ,,Magiczne przedmioty''.
Należy rzucić dwa razy 1K6.
Pierwszy rzut określa rodzaj magicznego przedmiotu, zaś drugi – jego dokładną nazwę, odczytywaną z kolumny, odpowiadającej ustalonemu pierwszym rzutem rodzajowi.

4.
Jeżeli postać posiada magiczne przedmioty należy ustalić ich rodzaj przy pomocy tabeli 14.9 ,,Magiczne przedmioty''.
Należy rzucić dwa razy 1K6.
Pierwszy rzut określa rodzaj magicznego przedmiotu, zaś drugi – jego dokładną nazwę, odczytywaną z kolumny, odpowiadającej ustalonemu pierwszym rzutem rodzajowi.

5.
Wszystkie rodzaje skarbów są dzielone między członków oddziału w sposób, odpowiadający wszystkim graczom.

Zasady szczegółowe:

[14.1] Skarby typu J, K i L zawsze są umieszczone w skrzyniach, które mogą być chronione przez pułapki.
Należy rzucić 1K6 dla każdej skrzyni – od 1 do 3 oczek oznacza, że w skrzyni założona jest pułapka.
Postać, posiadająca odpowiednie zdolności może próbować ją unieszkodliwić (patrz 7.0).
Jeżeli pułapka zostanie unieszkodliwiona lub jeżeli postacie przeżyją efekt jej działania dostęp do skarbów wewnątrz skrzyni stoi otworem.

[14.2] Rodzaj i wielkość skarbu należy ustalać dla każdego zabitego potwora osobno.

[14.3] Broń.

Kolumna ,,broń'' tabeli ,,Magiczne przedmioty'' reprezentuje rodzaj odebrane] potworowi magicznej broni.
Posiadanie przez postać danej broni musi być zaznaczone w jej karcie cech.

Po ustaleniu rodzaju broni należy rzucić trzeci raz kostką i przy pomocy tabeli 14.9 ,,Broń'' ustalić premię (tzn. liczbę, dodawaną do liczby oczek wyrzuconych walki) uzyskuje się ta broni, posługując się tą bronią.
Jeżeli odczytany z tabeli rezultat brzmi ,,rzuć dwa razy'' należy to zrobić, dodając do siebie uzyskane przy obu rzutach rezultaty.
Ponieważ istnieje możliwość, że wynik znów będzie równy ,,rzuć dwa razy'', więc premia, związana z posiadaniem magicznej broni teoretycznie może być nieskończenie duża.

Postać, posiadająca dwa rodzaje broni musi się jednego z nich pozbyć, jeżeli chce wziąć broń magiczną.

[14.4] Zbroja.

Noszenie magicznej zbroi wiąże się ze zwiększeniem siły postaci o taką liczbę punktów, jaką zostanie odczytana z tabeli ,,Magiczne przedmioty'' (kolumna ,,zbroja'') z rzędu, odpowiadającego liczbie wyrzuconych 1K6.
Postępowanie w przypadku uzyskania rezultatu ,,rzuć dwa razy'' jest takie samo, jak w punkcie 14.3.

[14.5] Mikstura.

Rodzaj mikstury ustala się rzucając 1K6 i odczytując odpowiadającą liczbie wyrzuconych oczek'' (kolumna ,,mikstura'') z tabeli ,,Magiczne przed-

mioty''.
Postacie muszą zbadać działanie mikstury, wypijając pewną jej ilość.
Jedna z postaci musi zostać wydelegowana do przeprowadzenia badania.
Postać ta natychmiast podlega działaniu mikstury.
Po przeprowadze-

niu próby zostaje tylko taka ilość mikstury, jaka wystarczy do jednorazowego użycia mikstury w przyszłości.
Działanie poszczególnych typów mikstury jest następujące:

Trucizna – Postać, przeprowadzająca badanie natychmiast odnosi 1K3 ran.

Zręczność – Zręczność postaci, która wypije tę miksturę jest zwiększona o 1K6 punktów, tylko w czasie najbliższej walki.
Potem mikstura traci moc.

Urok (osoba) – Postać, która wypiła tę miksturę może natychmiast zapanować w identyczny sposób, jak w wypadku czaru ,,Urok'' (patrz 10.5) nad Czarnym Rycerzem, Czarnoksiężnikiem, Orkiem lub Trollem, znajdującym się w tym samym segmencie.

Urok (potwór) – Działa tak samo jak mikstura urok (osoba), z tym, że daje władzę nad wszystkimi potworami, które nie są Czarnym Rycerzem, Czarnoksiężnikiem, Orkiem lub Trollem.

Uzdrowienie – Uzdrawia z 1K6 ran.

[14.6] Talizman.

Rodzaj talizmanu ustala się w ten sam sposób, jak wszystkich innych magicznych przedmiotów.
Jego działanie jest natychmiastowe.

Myśl – Postać, która go nosi może kosztem 1 rany zadać przeciwnikowi siłą swych myśli 1K6 ran (tylko wtedy, gdy atakowanemu w ten sposób potworowi nie uda się próba odparcia czaru).
Uwaga!
Ten talizman może zostać użyty tylko 1K6 razy.

Żółte słońce – Zwiększa potencjał magiczny, czyli jeżeli w tej przygodzie dominuje żółte słońce noszący ten talizman postać może używać tyle czarów więcej, ile oczek zostanie wyrzuconych 1K3.

Niebieskie słońce – j.w. dla niebieskiego słońca.

Czerwone słońce – j.w. dla czerwonego słońca.

Wszystkie słońca – Zwiększa o 1K3 potencjał magiczny dla każdego ze słońc.
Jeżeli w wyniku odnalezienia tego talizmanu (lub któregokolwiek z trzech wymienionych wyżej) postać może używać większej liczby czarów, niż ma zapisane w swej karcie cech gracz musi natychmiast dokonać wyboru nowych, dodatkowych czarów.

Zło – Postać, która weźmie do ręki ten talizman musi natychmiast podjąć próbę odparcia czaru.
Jeżeli próba się nie powiedzie postać zostaje opanowana przez zło i atakuje pozostałych członków oddziału.
Atakuje i zajmuje aktualnie zajmowaną w szyku pozycję.
Walka prowadzona jest aż: a. atakująca postać zostanie zabita lub b. pozostali członkowie oddziału zostaną zabici lub c. któryś z członków oddziału rzuci na atakującą postać Urok (za pomocą czaru lub mikstury) lub d. któryś z członków oddziału rzuci na nią czar ,,Uwolnienie''.

W przypadkach c. i d. opanowana przez zło postać musi starać się odeprzeć czar.
Niepowodzenie tej próby oznacza, że moc talizmanu została złamana i postać odzyskuje status normalnego członka oddziału.

[14.7] Medalion.

Rodzaj medalionu ustala się w ten sam sposób, jak wszystkich innych magicznych przedmiotów.
Jego działanie jest natychmiastowe.

Neutralizacja trucizny – Medalion daje postaci odporność na wszelkie trucizny tak